\chapter{Institutional Background and Hypothesis Development}
\label{chap:institutional_background}

This chapter establishes the institutional context and formal hypothesis structure underpinning the empirical analysis. It first explains the pre-liberalization structure of China's A-share market and the role of Stock Connect as implementation infrastructure for foreign participation. It then details the MSCI inclusion timeline and rule-based eligibility criteria, highlighting why the reform provides plausibly exogenous variation in international capital access. Finally, the chapter translates the conceptual framework from Chapter~\ref{chap:literature_review} into explicit testable hypotheses.

\section{China's Capital Market Structure Prior to Liberalization}

\subsection{The A-Share Market and Domestic Investor Segmentation}

China's A-share market---the renminbi-denominated equity market listed on the Shanghai and Shenzhen Stock Exchanges---was historically subject to strict foreign ownership restrictions and capital controls, which created near-complete segmentation from international capital flows~\cite{Cheng2024}. Before liberalization reforms, foreign investors could access mainland Chinese equities only through the limited QFII (Qualified Foreign Institutional Investor) and RQFII programs, which imposed strict quotas and administrative barriers. This segmentation meant that Chinese listed firms relied almost entirely on domestic capital markets, characterized by shorter investment horizons, weaker corporate governance standards, and limited exposure to global best practices.

\subsection{Stock Connect Programs as the Technical Infrastructure for Liberalization}

The technical infrastructure for capital market liberalization was established through the Stock Connect programs: Shanghai--Hong Kong Stock Connect (launched November 2014) and Shenzhen--Hong Kong Stock Connect (launched December 2016)~\cite{WangChia2014}. These programs developed daily and aggregate quotas, currency hedging methods, and settlement protocols that approved foreign investors to buy A-shares while adhering to capital control regimes. By 2018, these systems were sufficiently developed to handle the substantial capital movements associated with MSCI inclusion~\cite{Yin2023}.

\section{The MSCI Emerging Markets Index and the A-Share Inclusion Process}

\subsection{MSCI as a Global Benchmark: Scale, Influence, and Index Composition}

The MSCI Emerging Markets Index is the main benchmark for global institutional investors investing in developing market equities, with about \$1.9 trillion in assets under management tracking or comparing against MSCI emerging market indices in 2018. The index's composition determines mandatory portfolio allocations for passive funds and acts as a reference for active managers, influencing capital flows to the included companies~\cite{Wei2020}.

\subsection{Inclusion Criteria: Market Capitalization, Liquidity, and Stock Connect Accessibility}

MSCI's inclusion criteria were mechanical rather than discretionary: (i) a minimum free-float adjusted market capitalization threshold; (ii) minimum liquidity requirements based on annualized trading volume and turnover ratios; and (iii) accessibility through Stock Connect programs. Importantly, these criteria are independent of firms' underlying digitalization paths or performance, offering a plausible source of exogenous variation in international capital access.

\subsection{Pre-Announcement and Phased Implementation Timeline}

MSCI announced the inclusion decision on June 20, 2017---a full year before the first implementation tranche---providing firms with time to prepare for international investor engagement. The announcement itself constitutes a clean event that allows testing of announcement effects separately from implementation effects~\cite{Hu2018}.

\subsection{Phase~I (June 2018): 234 Large-Cap A-Shares at 5\% Inclusion Factor}

The initial tranche, launched on June 1, 2018, incorporated 234 large-cap A-shares at a 5\% inclusion factor. This accounted for about 0.73\% of the total MSCI Emerging Markets Index. The modest size aimed to evaluate market preparedness and reduce impact, while allowing the included companies to experience international passive fund inflows for the first time.

\subsection{Phase~II (November 2019): Expansion to 472 Securities at 20\% Factor}

By November 2019, the index grew to include 472 A-shares across large, mid, and small-cap categories, with a 20\% inclusion rate. This significantly boosted the presence of Chinese stocks in the index and broadened the scope to include smaller firms, expanding the universe of treated firms.

\subsection{Quasi-Experimental Properties of the Inclusion Design}

Three features make MSCI inclusion an almost perfect quasi-experiment. First, the selection criteria are mechanical and unrelated to digitalization propensity. Second, the timeline was externally imposed by MSCI, independent of Chinese government or firm decisions. Third, the phased implementation establishes clear pre- and post-treatment periods and multiple cohorts, enabling staggered adoption analysis.

\section{Hypothesis Development}

Building on the theoretical considerations and institutional background outlined above, we develop three testable hypotheses.

\subsection{H1: Main Effect --- Capital Market Liberalization and Digital Transformation Intensity}

\textbf{Hypothesis 1 (H1 --- Main Effect):} Capital market liberalization, as indicated by MSCI inclusion, is positively linked to the intensity of corporate digital transformation. Firms included in MSCI gain access to international capital, improved governance mechanisms, and global knowledge networks---all of which reduce the costs and increase the benefits of digital transformation investments.

\subsection{H2: Technology Heterogeneity --- Capital- and Knowledge-Intensive Domains}

\textbf{Hypothesis 2 (H2 --- Technology Heterogeneity):} The effects of capital market liberalization vary across digital technology types, with stronger effects for capital-intensive technologies (cloud computing, big data infrastructure) and knowledge-intensive technologies (artificial intelligence, advanced analytics) that particularly benefit from international financing and knowledge spillovers.

\subsection{H3: Regional and Firm Heterogeneity --- Absorptive Capacity and Institutional Quality}

\textbf{Hypothesis 3 (H3 --- Regional and Firm Heterogeneity):} Effects are moderated by regional development levels and firm-specific characteristics, with stronger treatment responses for firms in more developed regions with better institutional quality, larger firms with greater absorptive capacity, and firms in technology-intensive industries.
